\begin{lemma}[Hypersurfaces of maximal contact inside \(N\)]
\label{mc:contact-hypersurface-along}
\uses{mc:tangent-direction-along,mc:field-chain-rule,hir:order-along-flat,bex:hadamard,bcalc:snc-chart,bcalc:strict-ideal,bcalc:monomial-lift,bcalc:snc-simultaneous,ana:blowups,ana:local-holomorphic-inverse}
\source{wlodarczyk05-hironaka}
\lean{TranscendentalNakamaye.Pending.TangentDirectionAlong.exists_flatPartial_ne_zero,TranscendentalNakamaye.Pending.TangentDirectionAlong.contactCoords,TranscendentalNakamaye.Pending.TangentDirectionAlong.contactCoords_apply,TranscendentalNakamaye.Pending.TangentDirectionAlong.contactCoords_mem_plane_iff,TranscendentalNakamaye.Pending.TangentDirectionAlong.mdifferentiableAt_contactCoords,TranscendentalNakamaye.Pending.TangentDirectionAlong.injective_fderiv_contactCoords,TranscendentalNakamaye.Pending.TangentDirectionAlong.regularAtDim_inter_zeroSet,TranscendentalNakamaye.Pending.TangentDirectionAlong.exists_vanishingIdeal_inter_zeroSet_mem_iff,TranscendentalNakamaye.Pending.TangentDirectionAlong.IsTangentDirectionAlong.regularAtDim_contact,TranscendentalNakamaye.Pending.TangentDirectionAlong.isClosed_preimage_val_inter_zeroSet}
\leanok
Let \(\psi=\Phi\circ e_x\) be flattening coordinates of \(N\) at \(x\in N\), with \(N=\psi^{-1}(\mathbf C^k\times0)\) on their domain, let \(u\) be holomorphic on an open
\(U\ni x\) with \(u(x)=0\), and let \(\xi(u)(x)\ne0\) for a vector field \(\xi\) tangent to \(N\) at \(x\). Then
\begin{enumerate}
\item \(\partial u/\partial\psi_j(x)\ne0\) for some \(j<k\);
\item \(N\cap V(u)\cap U\) is regular of dimension \(k-1\) at \(x\), and near \(x\) its ideal is generated by \(u\) and the normal coordinates \(\psi_i\), \(i\ge k\):
\(\mathcal I_{N\cap V(u)}=\mathcal I_N+(u)\).
\end{enumerate}
Consequently, for a tangent direction \(u\) of \((\mathfrak a|_N,\mu)\) along a closed submanifold \(N\) on \(U\), \(N\cap V(u)\cap U\) is closed in \(U\) and regular, of
dimension \(\dim_yN-1\), at each of its points \(y\).
\end{lemma}
\begin{proof}
\leanok
(1) By Lemma~\ref{mc:field-chain-rule}, \(\xi(u)(x)=\sum_j\xi(\psi_j)(x)\,\partial u/\partial\psi_j(x)\); for \(j\ge k\), \(\psi_j\in\mathcal I_N\), so \(\xi(\psi_j)\in\mathcal I_N\) by
tangency and \(\xi(\psi_j)(x)=0\).
(2) Let \(l=k-1\) and \(\tau\) the transposition of \(j\) and \(l\). Put \(W_r=\psi_{\tau(r)}\) for \(r\ne l\) and \(W_l=u\). On the domain of \(\psi\), \(W\in\mathbf C^{k-1}\times0\)
if and only if \(u=0\) and \(\psi_r=0\) for \(r\ge k\) (\(\tau\) fixes these \(r\)), i.e.\ if and only if the point lies in \(N\cap V(u)\). In the coordinates \(\psi\), \(DW\) is
the selection matrix of \(\tau\) with the row \(l\) replaced by \(du\); a vector \(w\) in its kernel has \(w_m=0\) for \(m\ne j\), and then the row \(l\) gives
\(w_j\,\partial u/\partial\psi_j(x)=0\). So \(DW(x)\) is injective, and \(W-W(x)\), read in the chart at \(x\), is a local holomorphic coordinate system centred at \(x\)
(Lemma~\ref{bcalc:snc-chart}) in which \(N\cap V(u)\cap U\) is \(\mathbf C^{k-1}\times0\); since \(W(x)\in\mathbf C^{k-1}\times0\), \(W-W(x)\) lies in the plane exactly when \(W\) does.
Hadamard's lemma in these coordinates (Lemma~\ref{bcalc:strict-ideal}) generates the ideal by \(W_r\), \(r\ge k-1\), which are \(u\) and the \(\psi_r\), \(r\ge k\).
The consequence: at a point \(y\) of \(N\cap V(u)\cap U\) apply (2) with flattening coordinates of \(N\) at \(y\) and the field given by the tangent direction; closedness in
\(U\) holds because \(N\) is closed and \(u\) is continuous on \(U\).
\end{proof}